\documentclass[10pt, a4]{article}

% Packages:
\usepackage[
    ignoreheadfoot, % set margins without considering header and footer
    top=2 cm, % seperation between body and page edge from the top
    bottom=2 cm, % seperation between body and page edge from the bottom
    left=2 cm, % seperation between body and page edge from the left
    right=2 cm, % seperation between body and page edge from the right
    footskip=1.0 cm, % seperation between body and footer
    % showframe % for debugging 
]{geometry} % for adjusting page geometry
\usepackage{titlesec} % for customizing section titles
\usepackage{tabularx} % for making tables with fixed width columns
\usepackage{array} % tabularx requires this
\usepackage[dvipsnames]{xcolor} % for coloring text
\definecolor{primaryColor}{RGB}{0, 79, 144} % define primary color
\usepackage{enumitem} % for customizing lists
\usepackage{fontawesome5} % for using icons
\usepackage{amsmath} % for math
\usepackage[
    pdftitle={Arsene Mallet's CV},
    pdfauthor={Arsene Mallet},
    pdfcreator={LaTeX},
    colorlinks=true,
    urlcolor=primaryColor
]{hyperref} % for links, metadata and bookmarks
\usepackage[pscoord]{eso-pic} % for floating text on the page
\usepackage{calc} % for calculating lengths
\usepackage{bookmark} % for bookmarks
\usepackage{lastpage} % for getting the total number of pages
\usepackage{changepage} % for one column entries (adjustwidth environment)
\usepackage{paracol} % for two and three column entries
\usepackage{ifthen} % for conditional statements
\usepackage{needspace} % for avoiding page brake right after the section title
\usepackage{iftex} % check if engine is pdflatex, xetex or luatex

% Ensure that generate pdf is machine readable/ATS parsable:
\ifPDFTeX
    \input{glyphtounicode}
    \pdfgentounicode=1
    \usepackage[utf8]{inputenc}
    \usepackage{lmodern}
\fi

% Some settings:
\AtBeginEnvironment{adjustwidth}{\partopsep0pt} % remove space before adjustwidth environment
\pagestyle{empty} % no header or footer
\setcounter{secnumdepth}{0} % no section numbering
\setlength{\parindent}{0pt} % no indentation
\setlength{\topskip}{0pt} % no top skip
\setlength{\columnsep}{0cm} % set column seperation
\makeatletter
\let\ps@customFooterStyle\ps@plain % Copy the plain style to customFooterStyle
\patchcmd{\ps@customFooterStyle}{\thepage}{
    \color{gray}\textit{\small Arsène Mallet - Page \thepage{} sur \pageref*{LastPage}}
}{}{} % replace number by desired string
\makeatother
\pagestyle{customFooterStyle}

\titleformat{\section}{\needspace{4\baselineskip}\bfseries\large}{}{0pt}{}[\vspace{1pt}\titlerule]

\titlespacing{\section}{
    % left space:
    -1pt
}{
    % top space:
    0.3 cm
}{
    % bottom space:
    0.2 cm
} % section title spacing

\renewcommand\labelitemi{$\circ$} % custom bullet points
\newenvironment{highlights}{
    \begin{itemize}[
        topsep=0.10 cm,
        parsep=0.10 cm,
        partopsep=0pt,
        itemsep=0pt,
        leftmargin=0.4 cm + 10pt
    ]
}{
    \end{itemize}
} % new environment for highlights

\newenvironment{highlightsforbulletentries}{
    \begin{itemize}[
        topsep=0.10 cm,
        parsep=0.10 cm,
        partopsep=0pt,
        itemsep=0pt,
        leftmargin=10pt
    ]
}{
    \end{itemize}
} % new environment for highlights for bullet entries

\newenvironment{onecolentry}{
    \begin{adjustwidth}{
        0.2 cm + 0.00001 cm
    }{
        0.2 cm + 0.00001 cm
    }
}{
    \end{adjustwidth}
} % new environment for one column entries

\newenvironment{twocolentry}[2][]{
    \onecolentry
    \def\secondColumn{#2}
    \setcolumnwidth{\fill, 4.5 cm}
    \begin{paracol}{2}
}{
    \switchcolumn \raggedleft \secondColumn
    \end{paracol}
    \endonecolentry
} % new environment for two column entries

\newenvironment{header}{
    \setlength{\topsep}{0pt}\par\kern\topsep\centering\linespread{1.5}
}{
    \par\kern\topsep
} % new environment for the header

\newcommand{\placelastupdatedtext}{%
  \AddToShipoutPictureFG*{%
    \put(
        \LenToUnit{\paperwidth-2 cm-0.2 cm+0.05cm},
        \LenToUnit{\paperheight-1.0 cm}
    ){\vtop{{\null}\makebox[0pt][c]{
        \small\color{gray}\textit{Dernière mise à jour en avril 2026}\hspace{\widthof{Dernière mise à jour en septembre 2024}}
    }}}%
  }%
}%

% save the original href command in a new command:
\let\hrefWithoutArrow\href

% new command for external links:
\renewcommand{\href}[2]{\hrefWithoutArrow{#1}{\ifthenelse{\equal{#2}{}}{ }{#2 }\raisebox{.15ex}{\footnotesize \faExternalLink*}}}


\begin{document}
    \newcommand{\AND}{\unskip
        \cleaders\copy\ANDbox\hskip\wd\ANDbox
        \ignorespaces
    }
    \newsavebox\ANDbox
    \sbox\ANDbox{}

    \placelastupdatedtext
    \begin{header}
        \textbf{\fontsize{24 pt}{24 pt}\selectfont Arsène Mallet}

        \vspace{0.3 cm}

        \normalsize
        \mbox{{\color{black}\footnotesize\faMapMarker*}\hspace*{0.13cm}Palaiseau, France}%
        \kern 0.25 cm%
        \AND%
        \kern 0.25 cm%
        \mbox{\hrefWithoutArrow{mailto:arsene.mallet@telecom-paris.fr}{\color{black}{\footnotesize\faEnvelope[regular]}\hspace*{0.13cm}arsene.mallet@telecom-paris.fr}}%
        \kern 0.25 cm%
        \AND%
        \kern 0.25 cm%
        \mbox{\hrefWithoutArrow{tel:+33-782-86-56-66}{\color{black}{\footnotesize\faPhone*}\hspace*{0.13cm} +33 7 82 86 56 66}}%
        \kern 0.25 cm%
        \AND%
        \kern 0.25 cm%
        \mbox{\hrefWithoutArrow{https://linkedin.com/in/arsene-mallet}{\color{black}{\footnotesize\faLinkedinIn}\hspace*{0.13cm}arsene-mallet}}%
        \kern 0.25 cm%
        \AND%
        \kern 0.25 cm%
        \mbox{\hrefWithoutArrow{https://github.com/arsnm}{\color{black}{\footnotesize\faGithub}\hspace*{0.13cm}arsnm}}%
    \end{header}

    \vspace{0.3 cm - 0.3 cm}


    \section{Formation}

        \begin{twocolentry}{
        \textit{Palaiseau, France}

        \textit{Sept 2024 – Présent}}
            \textbf{Télécom Paris}

            \textit{Diplôme d'Ingénieur}
        \end{twocolentry}

        \vspace{0.10 cm}
        \begin{onecolentry}
            \begin{highlights}
                \item Spécialisation en Systèmes Embarqués et Data Science.
            \end{highlights}
        \end{onecolentry}

        \begin{twocolentry}{
        \textit{Lyon, France}

        \textit{Sept 2021 – Août 2024}}
            \textbf{Lycée La Martinière Monplaisir}

            \textit{Classes Préparatoires aux Grandes Écoles (CPGE)}
        \end{twocolentry}

        \vspace{0.10 cm}
        \begin{onecolentry}
            \begin{highlights}
                \item Filière MPSI/MP* (Mathématiques, Physique, Chimie, Informatique).
                \item Préparation intensive aux concours d'entrée des grandes
                    écoles d'ingénieurs françaises.
            \end{highlights}
        \end{onecolentry}


    \section{Projets Académiques \& Personnels}

        \begin{onecolentry}
            \textbf{Contrôleurs Matériels Personnalisés \& Traitement du Signal}
        \end{onecolentry}
        \vspace{0.05 cm}
        \begin{onecolentry}
            \begin{highlights}
                \item Conception et simulation de contrôleurs matériels en
                    VHDL/Verilog.
                \item Développement d'un contrôleur mémoire personnalisé
                    s'interfaçant avec de la Block RAM (BRAM), supportant les
                    modes de transaction simple et rafale (\textit{burst}) via
                    le bus Avalon.
                \item Architecture d'un contrôleur vidéo avec sortie VGA et
                    gestion de mémoire SDRAM, intégrant un arbitre matériel pour
                    coordonner l'accès aux données.
            \end{highlights}
        \end{onecolentry}
        \vspace{0.2 cm}

        \begin{onecolentry}
            \textbf{Programmation Embarquée Asynchrone en Rust}
        \end{onecolentry}
        \vspace{0.05 cm}
        \begin{onecolentry}
            \begin{highlights}
                \item Développement d'une application Rust \textit{bare-metal}
                    et multitâche, utilisant un framework asynchrone pour gérer
                    simultanément une machine virtuelle personnalisée, des
                    communications série et des tâches d'affichage en temps
                    réel.
                \item Conception d'une gestion de mémoire sans allocation
                    (\textit{zero-allocation}) sécurisée et d'une communication
                    inter-tâches (Channels, Signals, Mutexes) dans un
                    environnement \texttt{no\_std} strict.
            \end{highlights}
        \end{onecolentry}
        \vspace{0.2 cm}

        \begin{onecolentry}
            \textbf{Programmation Embarquée Bas Niveau (Bare-Metal)}
        \end{onecolentry}
        \vspace{0.05 cm}
        \begin{onecolentry}
            \begin{highlights}
                \item Développement d'un firmware C \textit{bare-metal}
                    utilisant une architecture \textit{super-loop} et des
                    interruptions matérielles (IRQ) pour gérer l'affichage de
                    données en temps réel sur une matrice de LED.
                \item Implémentation de pilotes (\textit{drivers}) bas niveau
                    incluant GPIO, Horloges, UART avec IRQ, I2C et Timers.
            \end{highlights}
        \end{onecolentry}
        \vspace{0.2 cm}

        \begin{onecolentry}
            \textbf{Simulation Hospitalière \& Prédiction par IA}
        \end{onecolentry}
        \vspace{0.05 cm}
        \begin{onecolentry}
            \begin{highlights}
                \item Système complet de simulation d'un hôpital intégrant une
                    simulation à événements discrets (SimPy), une surveillance
                    en temps réel (Flask/Chart.js) et des prédictions par
                    \textit{machine learning} (scikit-learn).
            \end{highlights}
        \end{onecolentry}
        \vspace{0.2 cm}

        \begin{onecolentry}
            \textbf{Système de Contrôle de Robot Autonome}
        \end{onecolentry}
        \vspace{0.05 cm}
        \begin{onecolentry}
            \begin{highlights}
                \item Localisation par vision (OpenCV), navigation autonome et
                    contrôle via interface web (Python, Raspberry Pi, contrôleur
                    I2C).
            \end{highlights}
        \end{onecolentry}
        \vspace{0.2 cm}

        \begin{onecolentry}
            \textbf{Simulation de Production en Usine}
        \end{onecolentry}
        \vspace{0.05 cm}
        \begin{onecolentry}
            \begin{highlights}
                \item Implémentation en Java suivant le patron de conception MVC
                    (Modèle-Vue-Contrôleur) avec interface graphique.
            \end{highlights}
        \end{onecolentry}


    \section{Expérience Professionnelle}

        \begin{twocolentry}{
        \textit{Lyon, France}

        \textit{2022 – 2026}}
            \textbf{Hospices Civils de Lyon (HCL)}

            \textit{Préparateur de Commandes \& Agent Logistique}
        \end{twocolentry}

        \vspace{0.10 cm}
        \begin{onecolentry}
            \begin{highlights}
                \item Travail à la Pharmacie Centrale des HCL lors des vacances
                    scolaires et estivales en tant que préparateur de commandes
                    et agent de logistique en magasin.
            \end{highlights}
        \end{onecolentry}

        \begin{twocolentry}{
        \textit{France}

        \textit{Sept 2024 – Présent}}
            \textbf{Complétude}

            \textit{Professeur Particulier - Mathématiques \& Physique}
        \end{twocolentry}

        \vspace{0.10 cm}
        \begin{onecolentry}
            \begin{highlights}
                \item Cours particuliers personnalisés pour des élèves de
                    collège et de lycée.
            \end{highlights}
        \end{onecolentry}


    \section{Compétences Techniques}

        \begin{onecolentry}
            \textbf{Langages de Programmation :} Python, C, Assembleur (ARM),
            SystemVerilog, Rust, Java, OCaml, C++, JavaScript.
        \end{onecolentry}

        \vspace{0.2 cm}

        \begin{onecolentry}
            \textbf{Outils \& Technologies :} Git, Linux, Programmation
            Embarquée (C/Rust \textit{bare-metal}, FPGA, QuestaSim), Data
            Science (Pandas, NumPy, Scikit-learn, PyTorch), SQL (PostgreSQL),
            Développement Web (HTML, CSS, JS, Astro), \LaTeX.
        \end{onecolentry}


    \section{Langues}

        \begin{onecolentry}
            \textbf{Français :} Langue maternelle.
        \end{onecolentry}

        \vspace{0.2 cm}

        \begin{onecolentry}
            \textbf{Anglais :} Courant (C1) - Certificat Linguaskill : 180+.
        \end{onecolentry}

        \vspace{0.2 cm}

        \begin{onecolentry}
            \textbf{Espagnol :} Intermédiaire (B1).
        \end{onecolentry}


    \section{Centres d'Intérêt}

        \begin{onecolentry}
            \textbf{Sports :} Pratique de divers sports, en particulier les
            sports de raquette et le cyclisme.
        \end{onecolentry}

        \vspace{0.2 cm}

        \begin{onecolentry}
            \textbf{Informatique \& Programmation :} Au-delà du cadre
            académique, j'aime comprendre le fonctionnement général des
            ordinateurs, coder sur mon temps libre et découvrir de nouvelles
            technologies.
        \end{onecolentry}

        \vspace{0.2 cm}

        \begin{onecolentry}
            \textbf{Jeux vidéo :} J'aime jouer aux jeux vidéo, particulièrement
            les jeux compétitifs.
        \end{onecolentry}


    \section{Informations Complémentaires}

        \begin{onecolentry}
            \textbf{Permis de conduire :} Permis B.
        \end{onecolentry}

        \vspace{0.2 cm}

        \begin{onecolentry}
            \textbf{Disponibilité :} À la recherche d'opportunités de stage dans
            les domaines de l'informatique, des systèmes embarqués et de la data
            science.
        \end{onecolentry}

\end{document}
